% ===========================================================================
\section{Related Work}
\label{sec:related}
% ===========================================================================

\subsection{Fuel-based and generation-based power inventories}

Power-sector emissions can be built from fuel consumed or from electricity generated with an
assumed heat rate. The IPCC 2006 Guidelines describe the fuel route as Tier~1 and treat
generation-based estimates as a fallback where fuel data are unavailable \citep{ipcc2006stationary},
because the heat rate is an efficiency assumption that varies by plant, by load and by year,
and because it is precisely the quantity a compiler is least able to verify.

Pakistan's regulator publishes fuel directly, so the fallback is unnecessary here. This is
worth stating because a workbook consolidating Pakistani plant-level data---which supplied the
captive register used in this paper---computes its \COtwo{} as kilograms per megawatt-hour
derived from IPCC factors through an assumed plant-efficiency range. That is the generation
route, and its output is not used here for that reason, while its plant register is.

\subsection{Coal quality and why the unit matters}

The two coals in Pakistan's fleet are imported bituminous, burned at Sahiwal, Port Qasim and
Hub, and domestic Thar lignite. Published proximate analyses of Thar coal report moisture near
\SI{46}{\percent}, ash around \SI{6.4}{\percent} and calorific values near
\SI{3030}{} kcal/kg---roughly half that of imported steam coal.

That difference is the whole issue. Emission factors expressed per tonne carry it in full;
factors expressed per terajoule do not, because the calorific value is exactly what the
differing energy content describes. The methodological literature on coal inventories treats
blend uncertainty as a major term, and it is---on a mass basis. Section~\ref{sec:methods}
shows it is a \SI{6.8}{\percent} term on an energy basis, which changes what a compiler should
worry about.

\subsection{Captive gas self-generation and the industrial-combustion gap}

Industrial self-generation sits awkwardly in inventory practice. IPCC assigns it to 1A2 with
manufacturing-industry factors, but the data that would support it are usually collected---if
at all---by an electricity regulator whose remit stops at plants selling into the grid.

For Pakistan the consequence is visible in the international comparisons: EDGAR
\citep{edgar2025booklet} reports Pakistani industrial combustion near \SI{44}{\mega\tonne} while most
national compilations report nothing, and the difference between a regulator's scope and an
inventory's scope is a substantial part of that. The activity data that do exist come from a
different regulator entirely---the oil and gas authority, which reports gas delivered to
captive power because it regulates the pipeline rather than the plant.

\subsection{Uncontrolled factors and control-device inventories}

US EPA's AP-42 compilation \citep{epaap42} publishes emission factors for coal, oil and gas
combustion predominantly as \emph{uncontrolled} values, with separate treatment of control
efficiencies, because what leaves a furnace and what leaves a stack differ by up to two orders
of magnitude for particulate and by most of the mass for sulphur.

Applying uncontrolled factors to a fleet with precipitators and desulphurisation, without
saying so, is a well-known way to overstate a power sector's air pollutants by an order of
magnitude. Section~\ref{sec:air} takes the position that where the control inventory does not
exist, the honest output is the uncontrolled bound explicitly labelled as one---not a
plausible-looking estimate resting on an invented removal efficiency.
